\documentclass[fleqn]{beamer}
\usepackage[english]{babel}

\usepackage[useregional]{datetime2}
\usepackage[utf8]{inputenc}
\usepackage{mathtools,amsthm,amssymb,bm,amsmath,mathrsfs}
\usepackage{amsfonts}
\usepackage{graphicx}
\usepackage[backend=biber,style=authoryear]{biblatex}
\addbibresource{biblio.bib}

\newcommand{\projdim}[1]{\mathbb{P}^#1}
\newcommand{\esperanza}[1]{\mathbb{E}(#1)}
\newcommand{\proba}[1]{\mathbb{P}(#1)}
\newcommand{\gavilla}[2]{\mathcal{O}_{#1}(#2)}

\newcommand{\citap}[1]{\textcolor{blue}{(\cite{#1})}}

% diagramas
\usepackage{tikz, tikz-cd}
\usepackage{pgf}
\usetikzlibrary{cd, babel}  % evita errores en tikz
\usetikzlibrary{graphs}
\usetikzlibrary{graphs.standard}
\usetikzlibrary{arrows,automata}


% vertical separator macro
\newcommand{\vsep}{
  \column{0.0\textwidth}
    \begin{tikzpicture}
      \draw[very thick,black!10] (0,0) -- (0,7.3);
    \end{tikzpicture}
}

% More space between lines in align
\setlength{\mathindent}{0pt}

% Beamer theme
\usetheme{ZMBZFMK}
\usefonttheme[onlysmall]{structurebold}
\mode<presentation>
%\setbeamercovered{transparent=10}

% align spacing
\setlength{\jot}{0pt}


\AtBeginSection[]
{
  \begin{frame}
    \frametitle{Table of Contents}
    \tableofcontents[currentsection]
  \end{frame}
}
\AtBeginSubsection[]
{
  \begin{frame}
    \frametitle{Table of Contents}
    \tableofcontents[currentsubsection]
  \end{frame}
}


% Only the first Slide
\title{Título}
\author{Horacio Castellanos Muñoa}
\institute[Universidad de Montevideo]{ }
\date{\today}


% Portada
\begin{document}
\begin{frame}
  \titlepage
\end{frame}

\section{Sección}
\subsection{Subsección}
\begin{frame}
  \frametitle{Título diapositiva}
  Esperanza $\esperanza{X}$. probabilidad $\proba{X}$ \begin{theorem} teorema
      
  \end{theorem}
  
  \begin{proof} Demostración \end{proof}

De acuerdo a \cite{dirac}, o en estilo natbib \citap{dirac}.
\end{frame}

\section{Gráficas}
\subsection{tikz}

\begin{frame}[fragile]{Como usar tikz}

\begin{tikzcd}
A \dar{g} \rar{f} & B \dar{h} \\
C \rar{i} & D
\end{tikzcd}
\\
Para que tikz funcione se deben definir los frames así: begin{frame}[fragile]{Como usar tikz}.

\end{frame}

\begin{frame}[fragile]{Otro ejemplo}
    \begin{center}
    

\begin{tikzcd}
      & Y \arrow[d, "f" description, dashed] \arrow[rd, "f_2"] \arrow[ld, "f_1"'] &       \\
X_{1} & X_{1} \times X_{2} \arrow[r, "\pi_2"'] \arrow[l, "\pi_2"]                 & X_{2}
\end{tikzcd}

\end{center}
\end{frame}

\section{Bibliografía}
\begin{frame}{Bibliografía}
\nocite{*}
\printbibliography
\end{frame}

\end{document}